\documentclass[11pt]{article}
\usepackage{tikz}
\usetikzlibrary{shapes, fit}
\usepackage{multirow}
\usepackage{booktabs}
\usepackage{rotating}
\usepackage{tabularx}
 \usepackage{xcolor,soul,framed} %,caption
\colorlet{shadecolor}{yellow}
% \usepackage{color,soul}
\usepackage{graphicx,epstopdf}
\usepackage{graphicx}
\graphicspath{{../pdf/}{../jpeg/}}
\DeclareGraphicsExtensions{.pdf,.jpeg,.png}
\DeclareGraphicsExtensions{.eps}
\usepackage{epstopdf}
% Removed invalid \epstopdfDeclareGraphicsRule
\usepackage[cmex10]{amsmath}
\usepackage{array}
\usepackage{amsmath}
\usepackage{mdwmath}
\usepackage{lipsum}
\usepackage{mdwtab}
\usepackage{eqparbox}
\usepackage{siunitx}
\usepackage{textcomp, gensymb}
\usepackage{amsfonts}
\usepackage[noadjust]{cite}
\usepackage[center]{caption}
\usepackage{subcaption}
\usepackage[font=small,skip=2pt]{caption}
\captionsetup{skip=0pt}
\usepackage{comment}
\usepackage{tikz}
\usetikzlibrary{shapes,fit,positioning}
\usetikzlibrary{shapes,arrows}
\usepackage{soul}
\usepackage{xcolor}
\usepackage{cite}
\sethlcolor{yellow} 
\usepackage{enumitem}
\usepackage{mathtools}
\usepackage{footnote}
\usepackage{url}
\def\UrlBreaks{\do\/\do-}
\usepackage{textcomp}
\usepackage{graphicx}
\graphicspath{{../pdf/}{../jpeg/}}
\DeclareGraphicsExtensions{.pdf,.jpeg,.png}
\DeclareGraphicsExtensions{.eps}
% Removed duplicate epstopdf and invalid conversion rule
\usepackage{multicol}
\usepackage{multirow}
\usepackage{rotating}
\usepackage{xcolor,soul,framed}
\usepackage{enumitem}
\usepackage{comment}
\usepackage[font=small,skip=0.4pt]{caption}
\captionsetup{skip=0pt}
\captionsetup{justification=centering}
\usepackage{subcaption}
\usepackage[noadjust]{cite}
% Removed invalid IEEEtran conditional
\usepackage[T1]{fontenc} 
\usepackage{amsmath}
\interdisplaylinepenalty=2500
% Removed newtxmath to avoid Bbbk conflict
\usepackage{bm}
\usepackage{titlesec}
\usepackage[margin=0.8in]{geometry}
\usepackage{amsmath,amssymb,booktabs,array,longtable}
\usepackage{hyperref}
\usepackage{enumitem}
\setlist{nosep}

\begin{document}

\begin{center}
{\Large\textbf{Shiv Nadar University Chennai}}\\
\textbf{CS3807 -- Deep Learning Laboratory}\\
\textbf{Experiment 3}\\
{\large\textbf{Implementation of Convolutional Neural Networks (CNNs) for Image Classification}}
\end{center}

\renewcommand{\arraystretch}{1.25}

\begin{tabular}{|p{3.8cm}|p{8cm}|p{2cm}|}
\hline
Degree \& Branch & B.Tech Artificial Intelligence \& Data Science & Semester V\\ \hline
Subject Code \& Name & CS3807 -- Deep Learning Laboratory & AY: 2026--27\\ \hline
\end{tabular}

\section*{1. Objective}

To understand the working principle of Convolutional Neural Networks by implementing convolution, pooling, feature map visualization, and image classification using TensorFlow/Keras.

\section*{2. Learning Outcomes}

Students will be able to

\begin{enumerate}
\item Explain the convolution operation.
\item Compute output dimensions using stride and padding.
\item Visualize feature maps.
\item Implement max pooling and average pooling.
\item Calculate CNN parameters.
\item Build a CNN for image classification.
\item Evaluate CNN performance.
\end{enumerate}

\section*{3. Background Theory}

A Convolutional Neural Network consists of

\begin{itemize}
\item Convolution Layer
\item Activation Layer
\item Pooling Layer
\item Fully Connected Layer
\end{itemize}

The convolution operation is

\[
Y(i,j)=\sum_m\sum_n X(i+m,j+n)K(m,n)
\]

where

\begin{itemize}
\item $X$ : Input Image
\item $K$ : Kernel (Filter)
\item $Y$ : Feature Map
\end{itemize}

The output feature map size is

\[
\frac{N-F+2P}{S}+1
\]

where

\begin{tabular}{ll}
$N$ & Input Size\\
$F$ & Filter Size\\
$P$ & Padding\\
$S$ & Stride
\end{tabular}
\subsection*{3.1 Common Convolution Kernels}

A \textbf{kernel} (or \textbf{filter}) is a small matrix that slides over an input image to extract important visual features such as edges, textures, corners, and smooth regions. During convolution, the kernel is multiplied element-wise with the corresponding image pixels, and the results are summed to produce a single value in the output feature map. Different kernels highlight different image characteristics.

\subsubsection*{1. Identity Kernel}

The identity kernel preserves the original image without modifying its pixel values.

\[
\begin{bmatrix}
0 & 0 & 0\\
0 & 1 & 0\\
0 & 0 & 0
\end{bmatrix}
\]

\textbf{Purpose:}
\begin{itemize}
\item Preserves the original image.
\item Used for understanding the convolution operation.
\end{itemize}

\subsubsection*{2. Sobel-X Kernel (Vertical Edge Detection)}

The Sobel-X kernel detects \textbf{vertical edges} by measuring intensity changes along the horizontal direction.

\[
\begin{bmatrix}
-1 & 0 & 1\\
-2 & 0 & 2\\
-1 & 0 & 1
\end{bmatrix}
\]

\textbf{Applications:}
\begin{itemize}
\item Object detection
\item Lane detection
\item Face recognition
\end{itemize}

\subsubsection*{3. Sobel-Y Kernel (Horizontal Edge Detection)}

The Sobel-Y kernel detects \textbf{horizontal edges} by measuring intensity changes along the vertical direction.

\[
\begin{bmatrix}
-1 & -2 & -1\\
0 & 0 & 0\\
1 & 2 & 1
\end{bmatrix}
\]

\textbf{Applications:}
\begin{itemize}
\item Text detection
\item Boundary extraction
\item Medical image analysis
\end{itemize}

\subsubsection*{4. Laplacian Kernel}

The Laplacian kernel detects edges in all directions using the second-order derivative of image intensity.

\[
\begin{bmatrix}
0 & -1 & 0\\
-1 & 4 & -1\\
0 & -1 & 0
\end{bmatrix}
\]

\textbf{Applications:}
\begin{itemize}
\item Edge enhancement
\item Satellite image analysis
\item Medical imaging
\end{itemize}

\subsubsection*{5. Sharpening Kernel}

This kernel enhances fine details by emphasizing high-frequency components.

\[
\begin{bmatrix}
0 & -1 & 0\\
-1 & 5 & -1\\
0 & -1 & 0
\end{bmatrix}
\]

\textbf{Applications:}
\begin{itemize}
\item Image enhancement
\item Optical Character Recognition (OCR)
\item Face recognition
\end{itemize}

\subsubsection*{6. Box Blur Kernel}

The Box Blur kernel smooths an image by replacing each pixel with the average of its neighboring pixels.

\[
\frac{1}{9}
\begin{bmatrix}
1 & 1 & 1\\
1 & 1 & 1\\
1 & 1 & 1
\end{bmatrix}
\]

\textbf{Applications:}
\begin{itemize}
\item Noise removal
\item Image smoothing
\item Image preprocessing
\end{itemize}

\subsubsection*{7. Gaussian Blur Kernel}

Gaussian blur smooths the image while preserving the overall structure better than a simple average filter.

\[
\frac{1}{16}
\begin{bmatrix}
1 & 2 & 1\\
2 & 4 & 2\\
1 & 2 & 1
\end{bmatrix}
\]

\textbf{Applications:}
\begin{itemize}
\item Noise reduction
\item Computer vision preprocessing
\item Feature extraction
\end{itemize}

\subsubsection*{8. Emboss Kernel}

The Emboss kernel creates a three-dimensional appearance by emphasizing directional changes in intensity.

\[
\begin{bmatrix}
-2 & -1 & 0\\
-1 & 1 & 1\\
0 & 1 & 2
\end{bmatrix}
\]

\textbf{Applications:}
\begin{itemize}
\item Texture analysis
\item Artistic image effects
\end{itemize}

\subsubsection*{9. Outline Detection Kernel}

This kernel highlights object boundaries by emphasizing regions with rapid intensity changes.

\[
\begin{bmatrix}
-1 & -1 & -1\\
-1 & 8 & -1\\
-1 & -1 & -1
\end{bmatrix}
\]

\textbf{Applications:}
\begin{itemize}
\item Image segmentation
\item Boundary detection
\item Shape extraction
\end{itemize}

\subsubsection*{10. Motion Blur Kernel}

This kernel simulates motion blur along the diagonal direction.

\[
\frac{1}{5}
\begin{bmatrix}
1 & 0 & 0 & 0 & 0\\
0 & 1 & 0 & 0 & 0\\
0 & 0 & 1 & 0 & 0\\
0 & 0 & 0 & 1 & 0\\
0 & 0 & 0 & 0 & 1
\end{bmatrix}
\]

\textbf{Applications:}
\begin{itemize}
\item Motion simulation
\item Image restoration
\item Video processing
\end{itemize}

\subsubsection*{Summary of Common Kernels}

\begin{center}
\renewcommand{\arraystretch}{1.3}
\begin{tabular}{|p{3cm}|p{2cm}|p{4cm}|p{4cm}|}
\hline
\textbf{Kernel} & \textbf{Size} & \textbf{Purpose} & \textbf{Applications}\\
\hline
Identity & $3\times3$ & Preserve original image & Baseline comparison\\
\hline
Sobel-X & $3\times3$ & Detect vertical edges & Object detection\\
\hline
Sobel-Y & $3\times3$ & Detect horizontal edges & Boundary detection\\
\hline
Laplacian & $3\times3$ & Detect edges in all directions & Medical imaging\\
\hline
Sharpen & $3\times3$ & Enhance details & Image enhancement\\
\hline
Box Blur & $3\times3$ & Smooth image & Noise removal\\
\hline
Gaussian Blur & $3\times3$ & Reduce noise & Computer vision\\
\hline
Emboss & $3\times3$ & 3D effect & Artistic filtering\\
\hline
Outline & $3\times3$ & Boundary extraction & Segmentation\\
\hline
Motion Blur & $5\times5$ & Simulate motion & Video processing\\
\hline
\end{tabular}
\end{center}
\section*{4. Numerical Example 1: Convolution}

Consider the following input image.

\[
X=
\begin{bmatrix}
1&2&3\\
4&5&6\\
7&8&9
\end{bmatrix}
\]

Kernel

\[
K=
\begin{bmatrix}
1&0\\
0&1
\end{bmatrix}
\]

Output

First position

\[
1(1)+2(0)+4(0)+5(1)=6
\]

Second position

\[
2(1)+3(0)+5(0)+6(1)=8
\]

Third position

\[
4(1)+5(0)+7(0)+8(1)=12
\]

Fourth position

\[
5(1)+6(0)+8(0)+9(1)=14
\]

Feature Map

\[
\begin{bmatrix}
6&8\\
12&14
\end{bmatrix}
\]

\section*{5. Numerical Example 2: Max Pooling}

Feature map

\[
\begin{bmatrix}
1&5&2&3\\
7&8&1&0\\
4&6&9&5\\
2&3&1&8
\end{bmatrix}
\]

Using a $2\times2$ max pooling filter,

Window 1

\[
\begin{bmatrix}
1&5\\
7&8
\end{bmatrix}
\rightarrow 8
\]

Window 2

\[
\begin{bmatrix}
2&3\\
1&0
\end{bmatrix}
\rightarrow3
\]

Window 3

\[
\begin{bmatrix}
4&6\\
2&3
\end{bmatrix}
\rightarrow6
\]

Window 4

\[
\begin{bmatrix}
9&5\\
1&8
\end{bmatrix}
\rightarrow9
\]

Output

\[
\begin{bmatrix}
8&3\\
6&9
\end{bmatrix}
\]

\section*{6. Numerical Example 3: Parameter Calculation}

Suppose

Input Image

\[
32\times32\times3
\]

Convolution Layer

\begin{itemize}
\item Filters = 16
\item Kernel = $3\times3$
\end{itemize}

Parameters

\[
(3\times3\times3+1)\times16
\]

\[
=(27+1)\times16
\]

\[
=448
\]

Therefore,

Total trainable parameters = 448.

\section*{7. Dataset}

Dataset: CIFAR-10

\begin{itemize}
\item Training Images : 50,000
\item Testing Images : 10,000
\item Classes : 10
\item Image Size : $32\times32\times3$
\end{itemize}

\section*{8. Experimental Procedure}

\textbf{Task 1}

Load CIFAR-10.

\begin{itemize}
\item Display ten sample images.
\item Print dataset dimensions.
\item Plot class distribution.
\end{itemize}

\bigskip

\textbf{Task 2}

Implement a convolution layer.

Compare

\begin{itemize}
\item Kernel = $3\times3$
\item Kernel = $5\times5$
\item Kernel = $7\times7$
\end{itemize}

Record feature map sizes.

\bigskip

\textbf{Task 3}

Study hyperparameters.

Compare

\begin{itemize}
\item Stride = 1
\item Stride = 2
\item Padding = Same
\item Padding = Valid
\end{itemize}

Compute output dimensions.

\bigskip

\textbf{Task 4}

Visualize feature maps after the first convolution layer.

Display at least eight feature maps.

\bigskip

\textbf{Task 5}

Compare

\begin{itemize}
\item Max Pooling
\item Average Pooling
\end{itemize}

Compare output size and accuracy.

\bigskip

\textbf{Task 6}

Construct the following CNN.

\[
Input
\rightarrow
Conv
\rightarrow
ReLU
\rightarrow
MaxPool
\rightarrow
Conv
\rightarrow
ReLU
\rightarrow
MaxPool
\rightarrow
Flatten
\rightarrow
Dense
\rightarrow
Softmax
\]

Train for

\begin{itemize}
\item Optimizer : Adam
\item Epochs : 20
\item Batch Size : 32
\end{itemize}

\bigskip

\textbf{Task 7}

Evaluate

\begin{itemize}
\item Accuracy
\item Precision
\item Recall
\item F1-score
\item Confusion Matrix
\item Classification Report
\end{itemize}

\section*{9. Mandatory Plots}

\begin{center}
    \includegraphics[width=0.85\textwidth]{sample_images-3.eps}
\end{center}
\textbf{Interpretation:} The sample images display the CIFAR-10 classes properly loaded and plotted in a grid. This confirms that the dataset parsing and decoding into $32\times32\times3$ RGB formats was successful.

\vspace{5mm}
\begin{center}
    \includegraphics[width=0.7\textwidth]{class_distribution-3.eps}
\end{center}
\textbf{Interpretation:} The bar chart demonstrates a perfectly uniform class distribution across all 10 categories. This balanced dataset prevents the model from developing a bias toward any majority class during training.

\vspace{5mm}
\begin{center}
    \includegraphics[width=0.75\textwidth]{feature_maps-3.png}
\end{center}
\textbf{Interpretation:} These feature maps from the first convolutional layer highlight the extracted spatial hierarchies. The filters have successfully identified primary shapes, edges, and color gradients from the input image.

\vspace{5mm}
\begin{center}
    \includegraphics[width=0.75\textwidth]{training_validation_accuracy-2.eps}
\end{center}
\textbf{Interpretation:} The training accuracy steadily climbs to $\sim87\%$ while the validation accuracy plateaus around $68\%$. This gap indicates that the model is beginning to overfit to the training data in later epochs.

\vspace{5mm}
\begin{center}
    \includegraphics[width=0.75\textwidth]{training_validation_loss-2.eps}
\end{center}
\textbf{Interpretation:} Corresponding with the accuracy graph, the training loss decreases continuously, but the validation loss begins to rise after epoch 10. This confirms the overfitting trend and suggests early stopping could be beneficial.

\vspace{5mm}
\begin{center}
    \includegraphics[width=0.85\textwidth]{confusion_matrix.png}
\end{center}
\textbf{Interpretation:} The confusion matrix shows strong true positive diagonal activations for classes like ship, automobile, and truck. However, the model struggles to differentiate between visually similar classes such as cats and dogs, leading to off-diagonal confusion.

\section*{10. Additional Exercises}

\begin{enumerate}
\item \textbf{Calculate the output size for a $64\times64$ image using a $5\times5$ kernel with stride 2 and padding 2.}
\\ Output = $\frac{64 - 5 + 2(2)}{2} + 1 = 32.5 \rightarrow 32$. The output size is $32\times32$.

\item \textbf{Compute the trainable parameters of a convolution layer having 64 filters of size $3\times3$ and an RGB input.}
\\ Parameters = $(3 \times 3 \times 3 + 1) \times 64 = 1792$.

\item \textbf{Compare ReLU and Sigmoid activation functions.}
\\ ReLU avoids the vanishing gradient problem and is computationally cheaper, making it ideal for hidden layers. Sigmoid squashes values between 0 and 1, making it better suited for binary output layers but prone to saturation.

\item \textbf{Replace Max Pooling with Average Pooling and compare the results.}
\\ Average pooling yielded slightly better validation stability but softer feature extraction, resulting in a marginal difference in final validation accuracy ($68.2\%$ vs $67.9\%$).

\item \textbf{Increase the number of convolution filters from 16 to 64 and analyze the effect on accuracy and computation time.}
\\ Increasing filters allows the network to learn more complex patterns, improving accuracy. However, it exponentially increases the parameter count, significantly raising computation time and memory usage.
\end{enumerate}

\section*{11. Results}

The Convolutional Neural Network was successfully trained and evaluated on the CIFAR-10 dataset. The results are tabulated below:

\vspace{5mm}
\begin{center}
\renewcommand{\arraystretch}{1.3}
\begin{tabular}{|p{5cm}|p{4cm}|}
\hline
\textbf{Metric} & \textbf{Value} \\
\hline
Training Accuracy & 87.01\% \\
\hline
Testing Accuracy & 68.00\% \\
\hline
Precision (Macro Avg) & 0.68 \\
\hline
Recall (Macro Avg) & 0.68 \\
\hline
F1-score (Macro Avg) & 0.68 \\
\hline
Number of Parameters & 167,562 \\
\hline
\end{tabular}
\end{center}

\subsection*{Dataset Dimensions}

\begin{center}
\renewcommand{\arraystretch}{1.3}
\begin{tabular}{|l|c|}
\hline
\textbf{Dataset Split} & \textbf{Dimensions} \\
\hline
Training Features (x\_train) & (50000, 32, 32, 3) \\
\hline
Training Labels (y\_train) & (50000, 1) \\
\hline
Testing Features (x\_test) & (10000, 32, 32, 3) \\
\hline
Testing Labels (y\_test) & (10000, 1) \\
\hline
\end{tabular}
\end{center}

\subsection*{Task 3: Output Dimensions}

\begin{center}
\renewcommand{\arraystretch}{1.3}
\begin{tabular}{|c|c|c|c|}
\hline
\textbf{Kernel} & \textbf{Stride} & \textbf{Padding} & \textbf{Output Shape} \\
\hline
$3\times3$ & 1 & Same & (None, 32, 32, 16) \\
\hline
$5\times5$ & 1 & Same & (None, 32, 32, 16) \\
\hline
$7\times7$ & 1 & Same & (None, 32, 32, 16) \\
\hline
$3\times3$ & 1 & Valid & (None, 30, 30, 16) \\
\hline
$3\times3$ & 2 & Same & (None, 16, 16, 16) \\
\hline
$3\times3$ & 2 & Valid & (None, 15, 15, 16) \\
\hline
\end{tabular}
\end{center}

\section*{12. Discussion}

\begin{enumerate}
\item \textbf{Why is convolution preferred over fully connected layers for images?}\\
Convolutional layers leverage parameter sharing and local spatial coherence. A fully connected layer would require millions of weights for a small image, destroying spatial relationships, whereas CNNs preserve the 2D structure efficiently.
\item \textbf{How does stride affect the feature map size?}\\
A larger stride moves the filter by more pixels at a time, skipping over overlapping regions. This aggressively downsamples the spatial dimensions of the feature map, reducing computational load.
\item \textbf{What is the role of padding?}\\
Padding adds artificial pixels (usually zeros) around the border of the input. It prevents the spatial dimensions from shrinking too quickly in deep networks and ensures that edge pixels are analyzed as thoroughly as central pixels.
\item \textbf{Why is pooling used?}\\
Pooling layers downsample the feature maps, making the representation smaller and more manageable. This introduces translation invariance, meaning the network can recognize features regardless of where they appear in the image.
\item \textbf{How do feature maps represent image characteristics?}\\
Early feature maps represent low-level features like straight edges, curves, and colors. As we go deeper into the network, feature maps combine these to represent high-level semantic shapes (e.g., tires, ears).
\item \textbf{Why do CNNs require fewer parameters than MLPs?}\\
Because of weight sharing. A single filter (e.g., $3\times3$) slides across the entire image, applying the exact same 9 weights everywhere. An MLP would require a unique weight for every single pixel connection.
\end{enumerate}

\section*{13. References}

\begin{enumerate}
\item Goodfellow et al., \emph{Deep Learning}.
\item Bishop, \emph{Pattern Recognition and Machine Learning}.
\item Haykin, \emph{Neural Networks and Learning Machines}.
\item TensorFlow Documentation.
\item CIFAR-10 Dataset Documentation.
\end{enumerate}

\end{document}